\documentclass[11pt]{article}

\usepackage[T1]{fontenc}
\usepackage{amsmath,amssymb}
\usepackage[margin=1in]{geometry}
\usepackage[hidelinks]{hyperref}

\newcommand{\F}{\mathbb{F}_{2^{128}}}
\newcommand{\eq}{\operatorname{eq}}

\title{Jagged Flock Commitments by Removing the Public Padding}
\author{}
\date{}

\begin{document}
\maketitle

\section{Conclusion}

The construction is straightforward once the public padding is removed
\emph{before} ring switching.

Let \(n\) be the number of actual BLAKE3 compressions and let
\[
    N=2^\tau\geq n
\]
be Flock's power-of-two batch size.  Flock may remain completely unchanged:
it proves a witness containing the \(n\) real compressions followed by \(N-n\)
copies of its fixed padding compression.  After Flock has reduced validity to
its final witness-evaluation claims, the verifier subtracts the publicly known
padding contribution.  The resulting claims concern the real witness followed
by zeros.  That zero-extended witness is what ring switching and the Jagged PCS
open.

Thus there is no need to change Flock's R1CS, zerocheck, lincheck, or
constant-wire pin.

\section{Split the witness into private and public parts}

Write the Boolean Flock witness as a function of an inner circuit index \(i\)
and an outer compression index \(j\):
\[
 z_{\mathrm{full}}(i,j)=
 \begin{cases}
   z_{\mathrm{real}}(i,j), & j<n,\\
   p(i),                   & j\geq n,
 \end{cases}
\]
where \(p\) is the fixed witness of the padding compression, namely the
compression used for hashing the all-zero block.

Let
\[
    S_n(j)=1[j\geq n]
\]
and define
\[
    z_0(i,j)
      =z_{\mathrm{full}}(i,j)+p(i)S_n(j).
\]
Addition is over characteristic two.  Therefore
\[
 z_0(i,j)=
 \begin{cases}
   z_{\mathrm{real}}(i,j), & j<n,\\
   0,                       & j\geq n.
 \end{cases}
\]

Flock proves claims about \(z_{\mathrm{full}}\).  We deterministically convert
each such claim into a claim about \(z_0\).

\section{The public correction factors}

Flock's final witness points separate into an inner point and an outer point.
The inner part uses Flock's univariate-skip convention, but it is still
tensor-separated from the outer instance coordinates.  Consequently,
\[
    \widehat z_0(x_{\mathrm{inner}},x_{\mathrm{outer}})
      =
    \widehat z_{\mathrm{full}}(x_{\mathrm{inner}},x_{\mathrm{outer}})
      +
    \widehat p(x_{\mathrm{inner}})
    \widehat S_n(x_{\mathrm{outer}}).
\]

Both factors in the correction are public:

\begin{itemize}
  \item \(\widehat S_n(x_{\mathrm{outer}})\) is a logarithmic-time suffix
        indicator evaluation;
  \item \(\widehat p(x_{\mathrm{inner}})\) is the evaluation of one fixed
        BLAKE3 padding witness.
\end{itemize}

The table has \(2^{14}=16{,}384\) Boolean positions.  For the current fixed
zero-block compression witness, exactly \(3{,}839\) positions are one; the
evaluation domain is therefore neither \(4{,}096\) nor \(8{,}192\) entries.

The padding witness is constant with respect to the outer compression
variables, which is what gives the factorization.  It is not constant with
respect to the inner circuit variables: evaluating it naively touches the
fixed \(2^{14}\)-bit witness.  This is nevertheless a constant cost independent
of \(n\) and \(N\), and the witness bits can be baked into the native and
Python verifiers.  The recursive verifier does not evaluate this table:
it defers the two values and merges all such claims with one 14-round
sumcheck, exactly like the existing bytecode and matrix deferrals.

The current Flock reduction produces two final witness claims, one for the
\(A\)-\(B\) side and one for the \(C\) side.  In general they have different
inner points, so the verifier performs this fixed-witness evaluation once per
claim.  The two evaluations can share fixed tables and implementation
machinery.

\section{Apply ring switching after the correction}

For every Flock output claim
\[
    \widehat z_{\mathrm{full}}(x)=v_{\mathrm{full}},
\]
derive, without receiving another prover value,
\[
    v_0
      =
    v_{\mathrm{full}}
      +
    \widehat p(x_{\mathrm{inner}})
    \widehat S_n(x_{\mathrm{outer}}).
\]
The claim passed to ring switching is then
\[
    \widehat z_0(x)=v_0.
\]

The prover likewise constructs ring switching's \(s_{\widehat v}\) data from
\(z_0\), or equivalently corrects the data obtained from
\(z_{\mathrm{full}}\) by the same public padding component.

Ring switching is linear, so no other change to the reduction is needed.  The
important point is that its input witness now has a zero suffix.

\section{Variable ordering and the Jagged prefix}

Choose the seven packed coordinates first, the remaining within-compression
coordinates next, and the compression-index coordinates last.  This is already
the natural current layout:
\[
    q_{\mathrm{pkd}}[128j+s],
    \qquad 0\leq s<128.
\]
One compression therefore occupies \(128\) consecutive elements of \(\F\).
After the public padding has been removed, the packed witness is
\[
    q_0=
      (q_{\mathrm{real}},
       0^{128(N-n)}).
\]
Its nonzero support is the single prefix
\[
    [0,128n).
\]

Equivalently, before packing, ring switching views the Boolean witness as
\(128\) slices of size \(2^{m-7}\).  With the instance coordinates ordered
last, every slice has the real data on the left and a zero suffix on the
right.  The same Jagged boundary applies to all slices, and packing them
pointwise gives the one \(q_{\mathrm{pkd}}\) prefix above.

The commitment stores only the first \(128n\) packed elements.  All following
ordinary Jagged blocks begin immediately after that prefix instead of after
the full \(128N\)-element Flock cube.

\section{Opening the committed prefix}

Ring switching produces the same transparent logical opening weight \(W\) as
before.  Since the omitted witness suffix is zero,
\[
    \langle q_0,W\rangle
      =
    \left\langle q_{\mathrm{real}},
      W|_{[0,128n)}\right\rangle.
\]
The target does not require another correction.  The prover inserts only the
first \(128n\) entries of the ring-switch basis into the stacked PCS opening.
The Jagged adapter evaluates the corresponding interval-restricted weight.

There is one standard Jagged subtlety: the verifier must evaluate the MLE of
the restricted weight.  It must not multiply two MLE evaluations, since
\[
 \operatorname{MLE}(W\cdot 1_{[0,128n)})
 \neq
 \operatorname{MLE}(W)\operatorname{MLE}(1_{[0,128n)})
\]
in general.

The existing ring-switch weight is a sum of \(128\) tensor terms.  The Jagged
interval evaluator can restrict each term and add the results.  Because the
interval starts at zero and is aligned to complete 128-element compression
rows, this specializes to the existing seven-coordinate product followed by a
small weighted prefix computation over the \(\tau\) instance coordinates.
Its cost remains \(O(128(7+\tau))\).

\section{Fixed BLAKE3 value slots}

The VM's virtual BLAKE3 value columns can be treated in either of two
equivalent ways:

\begin{enumerate}
  \item keep their current public padding values and subtract the corresponding
        one-slot suffix contribution before opening; or
  \item define their logical padding to be zero, matching \(q_0\), since these
        are artificial table-padding rows.
\end{enumerate}

The second choice gives the cleaner leanVM--Flock interface.  Actual table rows
are unchanged; only the public default for nonexistent rows becomes zero.
Their fixed-slot claims then become ordinary Jagged claims at logical point
\[
    (\operatorname{bits}(s),r)
\]
on the interval \([0,128n)\).

\section{Soundness}

The composition is sound provided that:

\begin{enumerate}
  \item the actual count \(n\) and the fixed padding witness \(p\) are bound by
        the statement before any relevant challenge;
  \item both Flock output claims are corrected using their exact inner and
        outer points;
  \item the prover's ring-switch data is formed from the same zero-extended
        witness \(z_0\);
  \item the PCS verifier evaluates the genuinely interval-restricted
        ring-switch weight;
  \item the native, recursive, and Python verifiers perform the identical
        public correction.
\end{enumerate}

Flock proves its original padded witness.  The public identity
\[
    z_0=z_{\mathrm{full}}+pS_n
\]
then converts its final claims into claims about a uniquely determined
zero-extended witness.  Ring switching and the Jagged PCS bind those converted
claims to the committed real prefix.  The prover receives no freedom in the
conversion.

\section{Cost and savings}

The native and Python verifiers perform a constant number of evaluations of
the fixed padding-compression witness, plus logarithmic suffix-indicator
evaluations.  This cost does not grow with the number of BLAKE3 instances.
Recursion exports those evaluations as deferred claims: a fresh transcript
randomly combines the two claims from every verified proof, and a single
14-round sumcheck reduces them to one evaluation of the fixed padding witness
for the outer native verifier.  Hence recursion pays only the sumcheck
verification and transparent terminal-weight costs, never a $2^{14}$-entry
in-circuit evaluation.

The physical Jagged area saved is exactly
\[
    128(N-n)
\]
elements of \(\F\), or
\[
    2048(N-n)\ \text{bytes}.
\]
When \(n\) is just above half a power of two, this removes almost half of the
Flock packed column.  As with the other Jagged columns, the largest commitment
and proof-size improvement occurs when the reduced total area lowers the
global dense PCS dimension.

\section{Implementation summary}

\begin{enumerate}
  \item Leave Flock's padded proof unchanged.
  \item Bake the fixed padding witness into the native and Python verifiers,
        and defer its recursive evaluations through one batched sumcheck.
  \item Correct Flock's two final witness claims by
        \(\widehat p(x_{\mathrm{inner}})
          \widehat S_n(x_{\mathrm{outer}})\).
  \item Produce ring-switch data for the corrected zero-suffix witness.
  \item Commit only \(q_{\mathrm{pkd}}[0..128n]\).
  \item Use Jagged interval weights for ring-switch and fixed-slot openings.
\end{enumerate}

The entire construction is captured by
\[
 \boxed{
   \text{Flock padded witness}
   =
   \text{committed zero-extended witness}
   +
   \text{public repeated padding}.
 }
\]

\end{document}
